\section{Discussion}\label{sec:discussion}

This section places the construction of Part~I relative to the dynamics of active inference and to existing participatory and structural positions in the foundations of physics. The dynamical statements of the first subsection are derived in Section~\ref{sec:mass} (status D); the positioning in the second and third subsections is interpretive (status S). Physical and phenomenal readings that rest on the postulates P1 to P10 are collected in Part~II.

\subsection{First-Order Dynamics: Stiffness, Friction, and the Kinetic-Metric Postulate}

The dynamics the construction defines is first order: the natural-gradient flow of Eq.~\eqref{eq:gauge_natural_gradient}, $\dot\xi = -G_\eta^{-1}\nabla\mathcal{F}$, where $G_\eta$ is the Fisher metric with each block divided by its rate ($\eta_\mu$ for the means, $\eta_\Sigma$ for the covariances). It has the form $\gamma\dot q = -\nabla V$ of the overdamped reduction of a mechanical particle $m\ddot q + \gamma\dot q = -\nabla V$, with the friction coefficient replaced by the tensor $G_\eta$ and no term in the inertial slot (Eq.~\eqref{eq:friction_form}). The comparison illustrates the regime only: the flow is a gradient flow on the statistical manifold and is not obtained as the overdamped limit of a Newtonian stochastic differential equation. First-order natural-gradient descent is also the regime of standard active inference, whose neuronal dynamics approximate natural-gradient descent on the variational free energy~\cite{DaCosta2021neural}, an intrinsic Riemannian gradient flow in the sense of Amari~\cite{Amari1998}; for cortical systems, where synaptic time constants dominate any putative inertial term, the overdamped regime is the appropriate one. Section~\ref{sec:mass} derives what precision does in this flow. The Hessian of $\mathcal{F}$ at consensus is a stiffness, so a high-precision agent has small compliance and a given force displaces it little. At an isolated equilibrium $q_i = p_i$ the stiffness equals the Fisher metric, so the linear relaxation rates are $\eta_\mu$ and $\eta_\Sigma$ whatever the precision. The linearized spectrum of the gradient flow is real, so there is no oscillation frequency to measure; complex eigenvalues appear only in the receiver-only truncation, whose field is not the gradient of $\mathcal{F}$. The only velocity-quadratic form of the dynamics is the Rayleigh dissipation of Eq.~\eqref{eq:rayleigh_dissipation}, and for a single rate its square root is the rate of the information clock (Eq.~\eqref{eq:clock_dissipation};~\cite[Ch.~7]{DennisGaugeVFE2026}).

The kinetic-metric postulate of Section~\ref{sec:velocity_quadratic} adds the inertia that the flow lacks, with its own second-order equation of motion, Eq.~\eqref{eq:second_order_postulate}, whose overdamped limit is the natural-gradient flow. It is adopted only under the reading of the variational principle as a candidate fundamental scaffolding, and the free-energy-principle canon does not imply it: Friston, Trujillo-Barreto and Daunizeau~\cite{Friston2008DEM} obtain second-order behavior through generalized coordinates of motion and a first-order variational filter, and Friston et al.~\cite{Friston2023PhysReports} decompose the flow into solenoidal and dissipative parts of a first-order stochastic vector field; neither writes a velocity-quadratic kinetic energy whose inertia tensor is the potential Hessian. Under the postulate as stated, with inertia $M_K$ equal to the stiffness $M$ at consensus, every undamped linear mode has unit frequency in flow-time units, so the scaling $\omega^2 \propto 1/\bar{M}_i$ is not a consequence of the postulate. A precision-dependent frequency requires an inertia tensor different from the stiffness tensor. Taking the bare prior precision $\bar{\Lambda}_p$ as inertia and the full Hessian $M = \bar{\Lambda}_p + K_{\mathrm{c}}$ as stiffness, with $K_{\mathrm{c}}$ the coupling contribution, gives $\omega^2 \in \operatorname{spec}(\bar{\Lambda}_p^{-1}M) = 1 + \operatorname{spec}(\bar{\Lambda}_p^{-1}K_{\mathrm{c}})$, which is a second and different postulate, testable and not an identity. On either postulate the Fisher information enters the dynamics the construction actually runs as a stiffness and a friction tensor, and belief configurations of large Fisher information are hard to displace, not hard to accelerate.

\subsection{Relation to Existing Participatory Approaches}

Of the strands of Wheeler's program in Table~\ref{tab:wheeler_correspondence}, the construction models W3 most directly: Wheeler's third question, how come ``one world'' out of many observer-participants, read with F{\o}llesdal's dictum that meaning is ``the joint product of all the evidence that is available to those who communicate''~\cite{Wheeler1990}. Agents compare laws only through the explicit transport $\Omega_{ij}$, the residual $D_{\mathrm{KL}}(q_i\|\Omega_{ij}[q_j])$ is the frame-independent part of their disagreement, and intersubjective agreement is a stationary point of the consensus free-energy functional of Section~\ref{sec:free_energy}; reading this as Wheeler's strand has status S. The companion theory describes the multiscale participation loop in the same register, as resembling Wheeler's proposal without being derived from it or evidence for it, and notes that none of its three typed forms of participation separates the program from a hierarchical model with learned priors, while formation could~\cite[Ch.~14]{DennisGaugeVFE2026}.

Rovelli's relational quantum mechanics~\cite{Rovelli1996}, an informational theory in Wheeler's spirit that relativizes states and values to arbitrary physical systems (Section~\ref{sec:related_work}), is the closest relative of the construction, which relativizes laws to agents and makes the comparison map between perspectives explicit. The construction is compatible with the relational reading without establishing that physical laws are observer-dependent, and its observer-dependence is one of representation and belief, not of fact. With Regime~I frames the cocycle closes (Lemma~\ref{thm:vanishing_holonomy}), a global frame exists, and every transported law lives in one classical probability space, so the construction satisfies the classical assumptions that Brukner's no-go theorem for observer-independent facts~\cite{Brukner2018} exploits and does not model the situation it addresses. The precise point of contact with the cross-perspective-links debate within relational quantum mechanics~\cite{AdlamRovelli2022} arises with independent links on the interaction graph (Section~\ref{sec:discrete_regime_ii}). There a nontrivial loop holonomy obstructs a global frame, and a normalized record whose law separates loop-holonomy classes exists at one base point, by the companion theory's theorem on an operational record sensitive to network-loop holonomy~\cite[Ch.~10]{DennisGaugeVFE2026}. This is a belief-level counterpart of the question whether the perspectives of different relational observers can be linked.

QBism~\cite{Fuchs2014, Fuchs2017} reads quantum probabilities as an agent's personalist degrees of belief and, as participatory realism, claims Wheeler's program by name (Section~\ref{sec:related_work}). The construction shares the personalist starting point, since each agent carries its own $q_i$, and in this respect, together with the relational comparison through transport, it is of the relational and QBist type. It departs from QBism in two places. It has no quantum content, and it adds a transport law between agents and a consensus functional whose stationary points model intersubjective agreement, whereas QBism keeps measurement outcomes personal to each agent~\cite{DeBrotaFuchsSchack2020} and leaves agreement to communication between agents. Read classically, the consensus mechanism is precision-weighted opinion pooling of the kind DeGroot analyzed~\cite{degroot1974}; its closest quantum-foundations relative is quantum Darwinism, in which objectivity arises from records redundantly available to many observers~\cite{Zurek2009}.

M{\"u}ller's observer-first program~\cite{Mueller2020} formalizes Wheeler's law without law~\cite{Wheeler1983} directly. Observer states are primitive, no laws are presupposed, and an external world with simple probabilistic laws, together with agreement between observers, emerges as an asymptotic statistical phenomenon through algorithmic probability. That program is more faithful to Wheeler than the present construction on law without law, on the refusal of a continuum, and on the primacy of registration; the present construction is more explicit about multi-agent transport, gauge structure, and consensus dynamics. Law without law is not modeled here. The functional, the belief family, and the divergence are fixed in advance, so their stationarity conditions are laws in the ordinary sense, and the $\mathrm{GL}(K_q)$ invariance of the divergence is already present for a single agent, a consistency requirement that consensus must respect (Section~\ref{sec:gauge_invariance_consensus} and Remark~\ref{rem:kretschmann}). A law-without-law version would need the functional itself to be selected by the population dynamics, which is open (status C).

The mechanism of Page and Wootters~\cite{PageWootters1983} derives dynamics from correlations between a clock subsystem and the rest of a stationary global state, and so addresses the timelessness behind Wheeler's ``no time''. The construction does not do this. Its histories are ordered by an algorithmic flow parameter, and the information clock of Section~\ref{sec:time_info} supplies duration along that order, not the order itself. A classical relational analog is available in the construction's own terms: parametrize the other agents' histories by one agent's Fisher length, which is strictly increasing wherever that agent's free-energy gradient is nonzero and constant at its stationary points~\cite[Ch.~7]{DennisGaugeVFE2026}, so that the relational clock fails exactly where that agent stops updating. This is a structural parallel (status S), not a derivation of time.

\subsection{Kantian and Structural-Realist Positioning}

This subsection places the author's interpretive stance, cognitive-first structural monism (Section~\ref{sec:cognitive_first}), relative to Kant and to structural realism; none of it is derived. The construction recasts Kant's distinction between phenomena and noumena~\cite{Kant1781} in a structural-realist register closest to Cassirer's neo-Kantian variant~\cite{Cassirer1923}. Worrall's epistemic structural realism~\cite{Worrall1989}, which suspends judgment about the nature of the relata, is a related precursor. The ontic structural realism of Ladyman~\cite{Ladyman1998} and of Ladyman and Ross~\cite{Ladyman2007} is a neighbor whose naturalism the position departs from: the cognitive-first order of derivation reverses the Primacy of Physics Constraint, a methodological principle of Ladyman and Ross's naturalized metaphysics, not the definition of ontic structural realism. On this reading information geometry on the fibers is the structure available to agents. The base manifold $\mathcal{C}$ is a presupposed smooth index manifold carrying no physical metric, called noumenal by analogy only. Along an agent's belief section a Riemannian metric on $\mathcal{C}$ is pulled back from the Fisher metric; computed with the covariant derivative of the agent's own connection it is unchanged under a passive change of frame~\cite[Ch.~B2]{DennisGaugeVFESuppB2026}, so differences between agents' pulled-back metrics reflect differences in their beliefs. What measurement protocols label distance and duration then correspond to induced metric intervals and accumulated Fisher length, and what they label inertia corresponds at most to the precision that acts as a stiffness in Section~\ref{sec:mass}. The frame-invariant data of the construction, namely the relative-entropy and evidence-bound values, the conjugacy classes of link holonomies, and the per-edge twist classes, are compatible with epistemic, moderate, and ontic structural realism alike, and nothing in the mathematics selects among these positions or the neo-Kantian reading~\cite[Ch.~14]{DennisGaugeVFE2026}. The treatment of beliefs and frames as agent-internal constructs, not structural representations of an external reality, is likewise compatible with the enactive reading of active inference advanced by Ramstead, Kirchhoff and Friston~\cite{Ramstead2019}, although that reading rests on an action-perception coupling that the present construction does not model. On the Kantian side the position disavows the apodictic a priori, forms fixed and valid for all time, and keeps a revisable constitutive role for the requirement that shareable descriptions be invariant under transport (Section~\ref{sec:gauge_invariance_consensus}), in the sense of the relativized a priori that Friedman~\cite{Friedman2001} develops from Reichenbach's 1920 analysis of relativity. Whether the within-agent pullback can carry a Lorentzian geometry is taken up under the postulates of Part~II (Section~\ref{sec:signature_resolution}), where an indefinite form is exhibited and its selection remains open.

\section{Summary of Part I}\label{sec:conclusion}

Part~I has constructed the classical, multi-agent, information-geometric objects on which the Wheeler correspondence of Table~\ref{tab:wheeler_correspondence} rests. Agents are local sections of associated bundles over a presupposed smooth index manifold $\mathcal{C}$ carrying no physical metric; each carries a belief $q_i$ on the state fiber and a generative model $s_i$ on the model fiber, expressed in its own frame, and agents compare laws only through the transport operators $\Omega_{ij}$ and $\tilde\Omega_{ij}$. At one base point the active symmetry is one global copy of the structure group and per-agent frames are passive re-expressions. Regime~I frames are removable by the relative-entropy invariance of Theorem~\ref{thm:glk_invariance}, whereas independent links on the interaction graph, with orientation reversal imposed, carry loop holonomy that no choice of frames removes. The consensus functional couples agents through transported relative entropies; its attention rows are Gibbs optima with reduced value $-\tau\log Z$, and intersubjective agreement is a stationary point of the consensus free-energy functional, which is the formal content of W3. The dynamics is a first-order natural-gradient flow in which precision is a stiffness and the Fisher metric a friction tensor (Section~\ref{sec:mass}).

Participation across scales (W4) rests on three constructions. The first is hierarchical coarse-graining. A meta-agent is the gauge-covariant M-projection barycenter of its constituents, the moment-matching minimizer of the weighted sum of $D_{\mathrm{KL}}(\Omega_{Ii}q_i\|q_I)$, and it obeys the covariance theorem (Theorem~\ref{thm:rg_covariance}). The coarse-graining map is an exact block-variable pushforward followed by a truncation to the Gaussian family, with no rescaling, coupling-space flow, or fixed point, and with blocking ratio $b > 1$ it admits at most $\log_b N$ steps on a network of $N$ agents (Section~\ref{sec:rg_status}). The second is the cross-scale prior of Eq.~\eqref{eq:cross_scale_shadow}, read across episodes: when the transition of an episode-indexed joint copies the transported parent latent into the child, the predictive prior of the companion theory's delayed-feedback proposition is exactly this formula, with the parent law filtered from the previous episode in the role of $q_I$~\cite[Ch.~12]{DennisGaugeVFE2026}. The same-time form used in the energy is an energy-level simplification (status S), not a prior of one joint. The third is a declared root (Section~\ref{sec:ouroboros}), a fixed top-level prior that closes the hierarchy as its boundary condition; every finite acyclic generative model of this type needs one~\cite[Ch.~12]{DennisGaugeVFE2026}, the belief-channel apex of that section is an energy-level closure and not a generative factor, and the root is the formal residue of Wheeler's ``no tower of turtles''.

What Part~I does not supply is listed under W5 in Table~\ref{tab:wheeler_correspondence}; in particular, its multiscale feedback relaxes to stationarity unless observations drive it and so is not a self-excited circuit. Part~II uses these objects to build the information clock (W1) and the pulled-back geometry (W2), adds a binary registration channel that gives the bit its registration meaning, and states the postulates P1 to P10 under which these objects receive physical readings. Whether the construction describes physical reality remains open.

\part*{Part II: The Speculative Program}\addcontentsline{toc}{part}{Part II: The Speculative Program}
